\PassOptionsToPackage{quiet}{fontspec}
\documentclass[UTF8]{beamer}
\usepackage{amsmath,amsthm,amssymb,mathrsfs}
\usepackage{fancyhdr}
\usepackage{ulem}
\usepackage{indentfirst}
\usepackage{latexsym}
\usepackage{amsfonts,amscd,mathtools,color}
\usepackage{bbm,dsfont,yfonts}
\usepackage[all,cmtip]{xy}
\input diagxy
\usepackage{hyperref}
\usepackage{verbatim,enumitem}
\usepackage{tikz-cd}
\usepackage{booktabs}
\usepackage{ctex}
\usetheme{Antibes}


\setbeamertemplate{footline}[page number] % To replace the footer line in all slides with a simple slide count uncomment this line

\setbeamertemplate{navigation symbols}{} % To remove the navigation symbols from the bottom of all slides uncomment this line

\theoremstyle{definition}
\newtheorem{thm}{定理}[section]
\newtheorem{lem}[thm]{引理}
\newtheorem{prop}[thm]{命题}
\newtheorem{cor}[thm]{推论}
\theoremstyle{definition}
\newtheorem{defn}[thm]{定义}
\newtheorem{ques}[thm]{Question}
\newtheorem{exmp}[thm]{例}
\newtheorem{rem}[thm]{注}

\newcommand\thua{\mathrel{\rotatebox[origin=c]{90}{$\twoheadrightarrow$}}}
\def\oto{{\bfig\morphism<180,0>[\mkern-4mu`\mkern-4mu;]\place(86,0)[\circ]\efig}}
\def\rto{{\bfig\morphism<180,0>[\mkern-4mu`\mkern-4mu;]\place(78,0)[\mapstochar]\efig}}

\DeclareMathOperator{\ldd}{/}
\DeclareMathOperator{\rdd}{\backslash}
\DeclareMathOperator{\with}{\&}
\DeclareMathOperator{\ver}{|}

\newcommand{\sub}{{\rm sub}}
\newcommand{\id}{{\rm id}}
\newcommand{\CQ}{\mathcal{Q}}
\newcommand{\CP}{\mathcal{P}}
\newcommand{\CS}{\mathcal{S}}
\newcommand{\sfe}{{\sf e}}
\newcommand{\sfm}{{\sf m}}
\newcommand{\sy}{{\sf y}}
\newcommand{\sQ}{{\sf Q}}
\newcommand{\setQ}{{\sf Set}(\sQ)}
\newcommand{\Set}{\sf Set}
\newcommand{\bv}{\bigvee}
\newcommand{\bw}{\bigwedge}
\newcommand{\SFQ}{\sQ\text{-}{\sf SemFil}}
\newcommand{\FQ}{\sQ\text{-}{\sf Fil}}
\newcommand{\Fil}{{\sf Fil}}
\newcommand{\QSup}{\sQ\text{-}{\sf Sup}} \newcommand{\lra}{\longrightarrow}
\newcommand{\ra}{\rightarrow}
\newcommand{\Ra}{\Rightarrow}
\newcommand{\lda}{\swarrow}
\newcommand{\rda}{\searrow}

\numberwithin{equation}{section}
\renewcommand{\theequation}{\thesection.\arabic{equation}}
\renewcommand{\baselinestretch}{1.3}
\parskip 5pt
\AtBeginSection[] {
\begin{frame}<beamer>{提纲}
\tableofcontents[currentsection]
\end{frame}}


\title{\textcolor{white}{模糊集Goguen范畴上的双反变幂集monad}}
\author{ 答辩人：鲁斯嘉}
\institute
{\begin{center}
\large{指导老师：张德学
\ \\

\ \\

\ \\
四川大学数学学院}
\end{center}
}
\date{2023年5月14日} 
\begin{document}	
\AtBeginSection[] {  \begin{frame}<beamer>{提纲}   \tableofcontents[currentsection]   \end{frame}}


\begin{frame}
\titlepage % Print the title page as the first slide
\end{frame}
\section{背景和预备知识}

\begin{frame}
1965年, %为了描述模糊现象, 
Zadeh引入了模糊集的概念. 从数学上说, 一个模糊集就是一个集合到单位区间 $[0,1]$ 的映射, 此时单位区间解释为逻辑值域. 这一概念很快被推广到了逻辑值域为完备格的情形. 
\end{frame}

\begin{frame}
为了给模糊集理论建立一个坚实的数学基础, Goguen 引入了范畴${\sf Set}(L)$,现在称为$L$-模糊集(简称$L$-集)的Goguen范畴, 这里$L$是一个完备格, 且经常带有额外的(逻辑)结构.
\end{frame}

\begin{frame}
\begin{defn}
设 $L$ 是完备格. 一个 $L$-模糊集指一个序对 $(X,\alpha)$, 其中 $X$ 是一个集合, $\alpha\colon X\lra L$ 是一个映射. 集 $X$ 称为 $L$-模糊集  $(X,\alpha)$ 的承载集, 完备格 $L$ 称为它的逻辑值域, 值 $\alpha(x)$ 称为 $x$ 的隶属度.%在明确承载集$X$和逻辑值域$L$时,可直接称$\alpha\colon X\to L$为一个模糊集.
两个$L$-模糊集之间的一个Goguen 态射$$f\colon(X,\alpha)\longrightarrow(Y,\beta)$$指一个映射$f\colon X\longrightarrow Y$,    满足Goguen条件$\alpha\le\beta\circ f.$ 

$L$-模糊集和Goguen 态射构成的范畴称为$L$-模糊集的Goguen范畴,记作${\sf Set}(L)$.
\end{defn}
{\footnotesize
J.A. Goguen, $L$-fuzzy sets, Journal of Mathematical Analysis and Applications 18 (1967) 145-174.}
\end{frame}

\begin{frame} {}

我们主要关心Goguen范畴 ${\sf Set}(L)$ 的范畴论性质, 特别是与幂对象有关的构造和性质.
\end{frame}

%\begin{frame}
%比如通过这个重要的性质，我们就知道集合范畴是它自身上的对偶monadic范畴。
% \begin{thm} 	反变幂集函子$\exp^{-1}\colon \sf Set^{\rm op}\lra\sf Set$是monadic函子.\end{thm}

%		{\footnotesize			E. Riehl, Category Theory in Context, Dover Publications, 2016.} 	\end{frame}

\begin{frame}
%我们也希望Goguen范畴上能有类似的结构来帮助我们更加深入的了解这个范畴。


集合范畴上的双反变幂集monad和它的submonad (特别是滤子monad和超滤monad) 在范畴论、代数学、拓扑学等学科中都有重要的作用. 

一个自然的问题:
范畴 ${\sf Set}(L)$上是否存在类似的构造?

\vskip 7pt


一个事实: 若 $L$ 不是单点集, 则Goguen范畴 ${\sf Set}(L)$不是 topos, 于是该范畴中幂对象的性质与集合范畴有很大的差异.
\end{frame}	

%\section{预备知识}
\begin{frame}
\begin{defn}
一个有单位元的quantale $\sQ=(\sQ,\with,k)$ 是一个以$\with$为乘法, $k$为单位的幺半群,并且满足以下条件:\begin{enumerate}[label={\rm(\roman*)}] \setlength{\itemsep}{0pt}
\item 集$\sQ$是一个完备格, 最大元记为$1$, 最小元记为$0$; 
\item 乘法$\with$对任意上确界分配,即对任意的$p,q,p_{i},q_{i}\in \sQ~(i\in I)$ 都有 		$$p\with\big(\bigvee\limits_{i\in I}q_{i}\big)=\bigvee\limits_{i\in I}p\with q_{i},\quad \big(\bigvee\limits_{i\in I}p_{i}\big)\with q=\bigvee\limits_{i\in I}p_{i}\with q,$$ 	其中$I$是一个指标集. \end{enumerate}
\end{defn}
{\footnotesize
K.I. Rosenthal, \itshape{Quantales and Their Applications}, \upshape Longman,   1990.}
\end{frame}

\begin{frame}
本文主要关心逻辑值域为一个有单位元的quantale时, Goguen范畴的范畴论性质.在无特殊说明的情况下, $\sQ$均表示一个有单位元的quantale $\sQ=(\sQ,\with,k)$.
\end{frame}

\begin{frame}
%一个一般的非交换的quantale的左乘和右乘得到的结果是不同的，它们的右伴我们分别称为左推出和右推出。
对于任意给定的$q\in \sQ$,映射$$-\ldd q\colon\sQ\longrightarrow \sQ,\quad r\ldd q=\bigvee\{p\in \sQ\ver p\with q\le r\}$$是$$-\with q\colon\sQ\longrightarrow \sQ$$的一个右伴,称为$\with$的左推出; 映射$$p\rdd -\colon\sQ\longrightarrow \sQ,\quad p\rdd r=\bigvee\{p\in \sQ\ver p\with q\le r\}$$是$$p\with-\colon\sQ\longrightarrow \sQ$$的一个右伴,称为$\with$的右推出. 
\end{frame}	

\begin{frame}
左右推出满足对任意的$p,q,r\in \sQ$,有$$p\le r\ldd q\Longleftrightarrow p\with q\le r\Longleftrightarrow q\le p\rdd r.$$进一步,若quantale $\sQ$是交换的(即对任意的$p,q\in \sQ$,有$p\with q=q\with p$),则有
$$p\ra q:=q\ldd p=p\rdd q,\quad\forall p,q\in\sQ.$$
\end{frame}

\begin{frame}
一些记号:\\
(1)设 $\sQ$ 是一个有单位元的quantale, $X$ 是一个集合.任给 $\lambda,\gamma\in\sQ^X$,  令$$\lambda\swarrow\gamma=\bigwedge\limits_{x\in X}\lambda(x)\ldd\gamma(x),$$	$$\gamma\searrow\lambda=\bigwedge\limits_{x\in X}\gamma(x)\rdd \lambda(x).$$  
\pause
%\vskip 3pt
(2)设$X,Y$是集合, $f\colon X\longrightarrow Y$是映射. 任给$\gamma\in \sQ^X$, $\gamma$在$f$下的像指模糊集  $$f(\gamma)\colon Y\to \sQ,\quad f(\gamma)(y)=\bigvee\{\gamma(x)\mid f(x)=y\}.$$   任给 $\lambda\in \sQ^{Y}$, $\lambda$在$f$下的原像指模糊集 $$f^{-1}(\lambda)\coloneqq \lambda\circ f .$$ 
\end{frame}

\section{Goguen范畴上的协变幂集monad}
\begin{frame}	2012年, 在 $\sQ$ 是交换的情形,  Eklund, Kortelainen, Stout等引入了	Goguen范畴 $\setQ$ 上的协变幂集monad. 他们的构造对非交换的情形也有效.

\vskip 11pt

{\footnotesize P. Eklund, J. Kortelainen, L.N. Stout, Adding fuzziness to terms and power objects using a monadic approach, Fuzzy Sets and Systems 192 (2012) 104-122.}
\end{frame}

\begin{frame}	任给$\setQ$的一个对象$(X,\alpha)$, 令 \[\mathscr{U}(X,\alpha)=(\sQ^X,\alpha^\downarrow),\] 其中对任意的 $\gamma\in\sQ^X$,  $$\alpha^\downarrow(\gamma)=\alpha\swarrow\gamma=\bw_{x\in X}\alpha(x)\ldd\gamma(x).$$   对于Goguen态射 $f\colon(X,\alpha)\lra(Y,\beta)$,有 \[\mathscr{U} f\colon(\sQ^X,\alpha^\downarrow)\lra (\sQ^Y,\beta^\downarrow), \quad \gamma\mapsto f(\gamma).\]  对应$f\mapsto \mathscr{U} f$ 定义了一个函子 \[\mathscr{U}\colon\setQ\lra\setQ,\] 称为 $\setQ$上的协变幂集函子.
\end{frame}

\begin{frame} 任给$\setQ$的对象$(X,\alpha)$, 映射
$$ m_{(X,\alpha)}\colon (\sQ^{\sQ^X},\alpha^{\downarrow\downarrow})\lra (\sQ^X,\alpha^\downarrow),\quad  m_{(X,\alpha)}(\land)= \bv_{\gamma\in\sQ^X}\land(\gamma)\with\gamma $$  和 \[ e_{(X,\alpha)}\colon (X,\alpha)\lra (\sQ^X,\alpha^\downarrow),\quad  e_{(X,\alpha)}(x) =k_x \] 都满足Goguen 条件, 由此可证$(\mathscr{U}, m, e)$是$\setQ$上的一个monad, 称为$\setQ$上的协变幂集monad. 
\end{frame}

%    \begin{frame}
%       \begin{prop}
%			$(\mathscr{U}, m, e)$的一个代数是一个序对 $((X,o),\alpha)$, 其中 $(X,o)$是一个余完备$\sQ$-格, $\alpha\colon X\lra\sQ$ 是$X$的一个模糊集, 并且 $$\sup\colon (\sQ^X,\alpha^\downarrow) \lra(X,\alpha),\quad \gamma\mapsto\sup\gamma$$满足Goguen 条件,即对任意的$\gamma\in\sQ^X$,有 \[ \alpha\swarrow\gamma\leq \alpha(\sup\gamma).\] 两个代数间的一个同态是一个映射 $f\colon X\lra Y$ 使得 $f\colon (X,\alpha)\lra(Y,\beta)$ 是 Goguen 态射并且 $f\colon(X,o_X)\lra(Y,o_Y)$ 是一个左伴.
%		\end{prop}
%    \end{frame}

\begin{frame}{}
\begin{prop}[$(\mathscr{U}, m, e)$的Eilenberg-Moore 代数] $(\mathscr{U}, m, e)$的一个代数是一个序对 $((X,\otimes),\alpha)$, 其中 $(X,\otimes)$ 是一个$\sQ$-模, $\alpha\colon X\lra\sQ$ 是$X$的一个模糊集,  满足以下条件: \begin{enumerate}[label={\rm(\roman*)}] \setlength{\itemsep}{0pt}
\item 任给$A\subseteq X$, 有$\bw_{x\in A}\alpha(x)\leq\alpha(\bv A)$;
\item 任给$r\in\sQ$ 和 $x\in X$,有 $\alpha(x)\ldd r\leq \alpha(r\otimes x)$.\end{enumerate}

两个代数之间的一个同态为一个映射$f\colon X\lra Y$ 使得$f\colon(X,\alpha)\lra(Y,\beta)$是Goguen 态射并且 $f\colon(X,\otimes_X)\lra(Y,\otimes_Y)$是 $\sQ$-模同态.\end{prop}
\end{frame}

\section{Goguen范畴上的双反变幂集monad}
%接下来我们介绍Goguen范畴上的双反变幂集monad，这是我们论文的主要成果。类似于集合范畴中的双反变幂集monad，Goguen范畴Set(Q)上的双反变幂集monad也是由一对反变函子构成的伴随生成的。
\begin{frame}
任给Goguen范畴$\sf Set(\sQ)$的一个对象$(X,\alpha)$,令$$\mathscr{P}(X,\alpha)=(\sQ^{X},\alpha^\uparrow),\quad\mathscr{P}^\dagger(X,\alpha)=(\sQ^X,\alpha_\dagger^\uparrow),$$其中对任意的$\gamma\in\sQ^X,$ $$\alpha^\uparrow(\gamma)=\gamma\swarrow\alpha,\quad\alpha_\dagger^\uparrow(\gamma)=\alpha\searrow\gamma.$$
任给Goguen 态射$f\colon (X,\alpha)\lra(Y,\beta),$ $$\mathscr{P}f\colon (\sQ^Y,\beta^\uparrow)\lra(\sQ^X,\alpha^\uparrow),\quad\lambda\mapsto\lambda\circ f$$与$$\mathscr{P}^\dagger f\colon (\sQ^Y,\beta^\uparrow_\dagger)\lra(\sQ^X,\alpha^\uparrow_\dagger),\quad\lambda\mapsto\lambda\circ f$$均满足Goguen条件.
\end{frame}

\begin{frame}
\begin{prop}函子
$\mathscr{P} \colon \sf Set(\sQ)^{\sf op}\lra\sf Set(\sQ)$是$\mathscr{P}^\dagger\colon \sf Set(\sQ)\lra\sf Set(\sQ)^{\sf op}$的右伴.
\end{prop}
\pause 
伴随 $\mathscr{P}^\dag \dashv\mathscr{P}$诱导了$\setQ$上的一个monad $$\mathfrak{P}=(\mathscr{P} \mathscr{P}^\dag,\mu,\eta),$$ 其中,  $\mu=\mathscr{P} \epsilon\mathscr{P}^\dag$, $\epsilon$是伴随 $\mathscr{P}^\dag \dashv\mathscr{P}$的余单位. Monad $\mathfrak{P}$ 称为$\setQ$上的双反变幂集monad.
\end{frame}

\begin{frame}
%这个定理给出了我们所构造的Goguen范畴上的协变幂集monad和双反变幂集monad的联系，这也从侧面说明了我们这样的构造是自然且合理的。
\begin{thm}当 $\sQ$是一个交换quantale时,协变幂集monad $(\mathscr{U},m,e)$是双反变幂集monad $(\mathfrak{P},\mu,\eta)$的一个submonad.
\end{thm}

\end{frame}


\begin{frame}{}
\begin{thm}[$(\mathfrak{P},\mu,\eta)$的Eilenberg-Moore 代数]\label{main1} $\mathscr{P}\colon \setQ^{\rm op}\lra\setQ $ 是monadic函子. \end{thm}
\end{frame}

%\section{Goguen范畴上的滤子monad}
%Monad P是双反变Q-幂集monad的一个提升这一事实是非常有用的。作为它的一个应用,我们说明了如何通过提升集合范畴中的Q-滤monad Q-Fil得到Goguen范畴上的双反变幂集monad P的一个submonad,前者是双反变Q-幂集monad \exp_Q^{-2}的一个submonad。同样的想法也可用于构建P的其他submonad。
\begin{frame} 接下来, 我们借助双反变幂集monad 构造了$\Set(\sQ)$上的滤子monad.
\begin{defn} \label{Q-filter} 集合$X$上的一个 $\sQ$-滤指的是一个映射$F\colon \sQ^X\lra \sQ$, 满足以下条件: 对任意的$\lambda,\gamma\in \sQ^X$ 以及$r\in\sQ$, \begin{enumerate}[label=(F\arabic*)] \setlength{\itemsep}{0pt}
\item \label{FF1} $F(k_X)\geq k$, 其中$k_X$ 是取值$k$的常值映射$X\lra\sQ$; \item \label{FF2} $F(\lambda)\wedge F(\gamma)\leq F(\lambda\wedge\gamma)$; \item \label{FF3} $\gamma\swarrow\lambda  \leq F(\gamma)\ldd F(\lambda);$
\item\label{FF4}   $F(r_X)\leq r$, 其中$r_X$是取值$r$的常值映射$X\lra\sQ$.   \end{enumerate}
\end{defn}

%由条件\ref{FF3} 不难验证, \ref{FF2}和\ref{FF4}中的不等号其实是等号.
\end{frame}

\begin{frame}
任给集合 $X$,把$X$上所有$\sQ$-滤 组成的集合记为$$\FQ(X),$$ 任给 $f\colon X\lra Y$ 和 $F\in\FQ(X)$,   $$f(F)\colon\sQ^Y\lra\sQ, \quad f(F)(\gamma)=F(\gamma\circ f) $$   是$Y$上的一个 $\sQ$-滤, 称为 $F$ 在 $f$ 下的像. %.由此,我们得到一个函子 \[\FQ\colon{\sf Set}\lra{\sf Set},\] 		称为集合范畴上的$\sQ$-滤函子.
\end{frame}

%	\begin{frame} 		 $\sQ$-滤函子 $\FQ$ 是双反变$\sQ$-幂集函子 $\exp_\sQ^{-2}\colon{\sf Set}\lra{\sf Set}$的一个子函子, 它可以做成 monad $(\exp_\sQ^{-2},\mu, \eta)$的一个submonad $(\FQ,\mu,\eta)$. 	\end{frame}

\begin{frame}
任给 $\setQ$的对象$(X,\alpha)$, 令 \[\mathfrak{F}(X,\alpha)= (\FQ(X),(\alpha_\dag^\uparrow)^\uparrow),\] 其中, 对$X$上的每个  $\sQ$-滤 $F$,   \[(\alpha_\dag^\uparrow)^\uparrow(F) = \bw_{\gamma\in\sQ^X}F(\gamma)\ldd(\alpha\searrow\gamma).\] 		由此我们得到一个函子 $$\mathfrak{F}\colon\setQ\lra\setQ,$$  它是函子双反变幂集函子$\mathfrak{P}$的一个子函子.  \end{frame}

\begin{frame}
双反变幂集monad $(\mathfrak{P},\mu,\eta)$的乘法$\mu$对$\mathfrak{F}$封闭,单位$\eta$可经过$\mathfrak{F}$分解, 于是

\begin{prop} 三元组\((\mathfrak{F},\mu,\eta)\)是双反变幂集 monad $(\mathfrak{P},\mu,\eta)$ 的一个submonad. \end{prop}
\end{frame}

\section*{总结}
\begin{frame}{总结}
主要结果：
\begin{enumerate}[label={(\arabic*)}] \setlength{\itemsep}{0pt}
\item  当$\sQ$是交换quantale时, $\Set(\sQ)$上的协变幂集monad是双反变幂集monad $(\mathfrak{P},\mu,\eta)$ 的一个submonad.

\item 双反变幂集monad $(\mathfrak{P},\mu,\eta)$的 Eilenberg-Moore 代数范畴与$\Set(\sQ)$的对偶范畴等价, 于是, 与集合范畴类似, $\Set(\sQ)$也是它自身上的对偶monadic范畴.

\item  借助双反变幂集monad, 我们构造了$\Set(\sQ)$上的滤子monad.   \end{enumerate}
\end{frame}

\section*{致谢}
\begin{frame}

\centering\Huge 谢\ \ 谢！
\end{frame}

\end{document}
